\subsection{M. Riesz 与 Kolmogorov 的定理}
\label{sec:ch3-riesz-kolmogorov-theorems}

下一个定理说明, 目前定义在 $\mathcal S$ 或 $L^2$ 上的 Hilbert 变换可以延拓到 $L^p$ 中的函数, $1\le p<\infty$.

\begin{Theorem}
\label{thm:ch3-riesz-kolmogorov}
对 $f\in\mathcal S(\mathbb R)$, 下列断言成立:

\begin{enumerate}
\item (Kolmogorov) $H$ 为弱 $(1,1)$ 型:
\[
 \bigl|\{x\in\mathbb R:|Hf(x)|>\lambda\}\bigr|
 \le \frac C\lambda\|f\|_1.
\]

\item (M. Riesz) $H$ 为强 $(p,p)$ 型, $1<p<\infty$:
\[
 \|Hf\|_p\le C_p\|f\|_p.
\]
\end{enumerate}
\end{Theorem}

\hypertarget{chapter-03-page-064}{}
\begin{Proof}
(1) 固定 $\lambda>0$ 并设 $f$ 非负. 在高度 $\lambda$ 对 $f$ 作 Calderón--Zygmund 分解(见定理~\ref{thm:ch2-calderon-zygmund-decomposition}), 得到一列互不相交的区间 $\{I_j\}$, 使得
\[
 f(x)\le\lambda\quad\text{a.e. }x\in\Omega^c,
 \qquad \Omega=\bigcup_j I_j,
\]
\[
 |\Omega|\le\frac1\lambda\|f\|_1,
 \qquad
 \lambda<\frac1{|I_j|}\int_{I_j}f\le2\lambda.
\]
由这个 $\mathbb R$ 的分解, 将 $f$ 写成两个函数 $g$ 与 $b$ 之和, 其中
\[
 g(x)=
 \begin{cases}
  f(x),&x\notin\Omega,\\
  \displaystyle\frac1{|I_j|}\int_{I_j}f,&x\in I_j,
 \end{cases}
 \qquad
 b(x)=\sum_j b_j(x),
\]
并且
\[
 b_j(x)=\left(f(x)-\frac1{|I_j|}\int_{I_j}f\right)\chi_{I_j}(x).
\]
于是 $g(x)\le2\lambda$ 几乎处处成立, 而 $b_j$ 支撑于 $I_j$ 且积分为零. 因为 $Hf=Hg+Hb$,
\[
\begin{aligned}
 \bigl|\{x\in\mathbb R:|Hf(x)|>\lambda\}\bigr|
 &\le \bigl|\{x\in\mathbb R:|Hg(x)|>\lambda/2\}\bigr|\\
 &\quad+\bigl|\{x\in\mathbb R:|Hb(x)|>\lambda/2\}\bigr|.
\end{aligned}
\]
利用\eqref{eq:ch3-hilbert-l2-isometry} 估计第一项:
\[
\begin{aligned}
 \bigl|\{x\in\mathbb R:|Hg(x)|>\lambda/2\}\bigr|
 &\le\left(\frac2\lambda\right)^2\int_{\mathbb R}|Hg(x)|^2\,dx\\
 &=\frac4{\lambda^2}\int_{\mathbb R}g(x)^2\,dx
 \le\frac8\lambda\int_{\mathbb R}g(x)\,dx
 =\frac8\lambda\int_{\mathbb R}f(x)\,dx.
\end{aligned}
\]

\hypertarget{chapter-03-page-065}{}
令 $2I_j$ 表示与 $I_j$ 同心而长度为其两倍的区间, 并令 $\Omega^*=\bigcup_j2I_j$. 则 $|\Omega^*|\le2|\Omega|$, 而且
\[
\begin{aligned}
 \bigl|\{x\in\mathbb R:|Hb(x)|>\lambda/2\}\bigr|
 &\le|\Omega^*|+\bigl|\{x\notin\Omega^*:|Hb(x)|>\lambda/2\}\bigr|\\
 &\le\frac2\lambda\|f\|_1+\frac2\lambda\int_{\mathbb R\setminus\Omega^*}|Hb(x)|\,dx.
\end{aligned}
\]
几乎处处有 $|Hb(x)|\le\sum_j|Hb_j(x)|$: 当求和有限时这是显然的; 否则, 它来自 $\sum b_j$ 与 $\sum Hb_j$ 分别在 $L^2$ 中收敛到 $b$ 与 $Hb$. 因而要完成弱 $(1,1)$ 不等式的证明, 只需证明
\[
 \sum_j\int_{\mathbb R\setminus2I_j}|Hb_j(x)|\,dx\le C\|f\|_1.
\]
尽管 $b_j\notin\mathcal S$, 当 $x\notin2I_j$ 时公式
\[
 Hb_j(x)=\int_{I_j}\frac{b_j(y)}{x-y}\,dy
\]
仍成立. 记 $I_j$ 的中心为 $c_j$. 因为 $b_j$ 的积分为零,
\[
\begin{aligned}
 \int_{\mathbb R\setminus2I_j}|Hb_j(x)|\,dx
 &=\int_{\mathbb R\setminus2I_j}
   \left|\int_{I_j}b_j(y)\left(\frac1{x-y}-\frac1{x-c_j}\right)dy\right|dx\\
 &\le\int_{I_j}|b_j(y)|
   \left(\int_{\mathbb R\setminus2I_j}
   \frac{|y-c_j|}{|x-y|\,|x-c_j|}\,dx\right)dy\\
 &\le\int_{I_j}|b_j(y)|
   \left(|I_j|\int_{\mathbb R\setminus2I_j}
   \frac{dx}{|x-c_j|^2}\right)dy.
\end{aligned}
\]
最后一个不等式利用了 $|y-c_j|\le|I_j|/2$ 以及 $|x-y|\ge|x-c_j|/2$. 括号内的积分等于 $2$, 因此
\[
 \sum_j\int_{\mathbb R\setminus2I_j}|Hb_j(x)|\,dx
 \le2\sum_j\int_{I_j}|b_j(y)|\,dy\le4\|f\|_1.
\]
上面的弱 $(1,1)$ 证明针对非负函数 $f$, 但这已经足够, 因为任意实函数可分解为正部与负部, 复函数可分解为实部与虚部.

\hypertarget{chapter-03-page-066}{}
(2) 因为 $H$ 为弱 $(1,1)$ 型和强 $(2,2)$ 型, Marcinkiewicz 插值定理给出 $1<p<2$ 时的强 $(p,p)$ 不等式. 当 $p>2$ 时, 应用\eqref{eq:ch3-hilbert-skew-adjoint} 及 $p'<2$ 时的结果:
\[
\begin{aligned}
 \|Hf\|_p
 &=\sup\left\{\left|\int Hf\cdot g\right|:\|g\|_{p'}\le1\right\}\\
 &=\sup\left\{\left|\int f\cdot Hg\right|:\|g\|_{p'}\le1\right\}\\
 &\le\|f\|_p\sup\left\{\|Hg\|_{p'}:\|g\|_{p'}\le1\right\}
 \le C_{p'}\|f\|_p.
\end{aligned}
\]
\end{Proof}

定理~\ref{thm:ch3-riesz-kolmogorov} 第一部分证明中的函数 $g$ 与 $b$ 传统上分别称为 $f$ 的好部与坏部.

从这些证明还可看出, 强 $(p,p)$ 与 $(p',p')$ 不等式中的常数相同, 并且当 $p$ 趋于 $1$ 或无穷时趋于无穷. 更准确地说,
\[
 C_p=O(p)\quad(p\to\infty),
 \qquad
 C_p=O\bigl((p-1)^{-1}\bigr)\quad(p\to1).
\]

利用定理~\ref{thm:ch3-riesz-kolmogorov} 中的不等式, 可将 Hilbert 变换延拓到 $L^p$, $1\le p<\infty$. 若 $f\in L^1$, 且 $\{f_n\}$ 是 $\mathcal S$ 中依 $L^1$ 范数收敛到 $f$ 的函数列, 则由弱 $(1,1)$ 不等式, $\{Hf_n\}$ 在测度意义下为 Cauchy 列: 对任意 $\varepsilon>0$,
\[
 \lim_{m,n\to\infty}
 \bigl|\{x\in\mathbb R:|(Hf_n-Hf_m)(x)|>\varepsilon\}\bigr|=0.
\]
因此它在测度意义下收敛到一个可测函数, 我们将此函数定义为 $f$ 的 Hilbert 变换.

若 $f\in L^p$, $1<p<\infty$, 且 $\{f_n\}$ 是 $\mathcal S$ 中依 $L^p$ 范数收敛到 $f$ 的函数列, 则由强 $(p,p)$ 不等式, $\{Hf_n\}$ 是 $L^p$ 中的 Cauchy 列, 因而收敛到 $L^p$ 中的一个函数, 我们称之为 $f$ 的 Hilbert 变换.

在两种情形下, $\{Hf_n\}$ 都有一个依赖于 $f$ 的子列几乎处处逐点收敛到上述定义的 $Hf$.

当 $p=1$ 或 $p=\infty$ 时, 强 $(p,p)$ 不等式不成立. 取 $f=\chi_{[0,1]}$ 即可看出这一点. 此时
\[
 Hf(x)=\frac1\pi\log\left|\frac{x}{x-1}\right|,
\]
而 $Hf$ 既不可积也无界.

\hypertarget{chapter-03-page-067}{}
在 Schwartz 类中, 容易刻画 Hilbert 变换可积的函数: 对 $\varphi\in\mathcal S$, $H\varphi\in L^1$ 当且仅当 $\int\varphi=0$. 其证明留作练习. 关于 Hilbert 变换在 $L^1$ 上性质的更多资料, 见第~\ref{subsec:ch3-hilbert-l1}~节.

在\MBUnresolvedReference{chapter}{第 6 章}中, 我们将考察 Hilbert 变换仍然可积的可积函数空间, 以及能够为有界函数 $f$ 定义 $Hf$ 的一个比 $L^\infty$ 更大的空间(另见第~\ref{subsec:ch3-l-log-l-estimates}~节).
